\chapter{Larios Scale 2 Results}

Larios et al. 2018 present numerical evidence for blow-up of the Euler-Voigt equations from (\ref{eq:TGV}) occurring at $T^* \approx 4$. Their numerical model used parameters $L=1$ and $V_0=1$, and so the characteristic time-scale of their model is $t_c = 1$. These parameters give their model an initial seminorm $\|\bn\eta_{TGV}\|_{L^2} = \sqrt{3}\pi$ (\ref{app:tg:seminorm:1}) and initial vorticity $\|\nabla\times\eta_{TGV}\|_{L^\infty} = 4\pi$ (\ref{app:tg:vorticity:1}). 
\\~\\
Our numerical models are performed on a domain with the same size and so does not need to be rescaled. However, at this scaling the results of our model do not match the results of Fig 3. Therefore, we introduce an ad hoc scaling to $\|\bn\eta_{TGV}\|_{L^2} = \pi$.
\\~\\
The plots below show our results for the scaled TGV. In addition they show the comparison of our results with those of (Fig. 1b, Bustamante et al. 2012), (Fig. 2, Larios et al. 2018), and (Fig. 3, Larios et al. 2018). 

\section{Evolution of $\|\nabla\ua\|_{L^2}$}

This plot shows the evolution of the $\|\nabla\ua\|_{L^2}$ norm over time. Neither Bustamante, Brachet, or Larios show an equivalent plot but it is still worth examining. The growth in the norm strictly increases over time. Additionally, the norm at T increases as $\alpha \rightarrow 0$ which is consistent with the expected behaviour. The time-scale of the results is between the Bustamante scaled model and the first Larios scaling.

\begin{figure}[h!]
\centering
\includegraphics[scale=0.75]{"./Figures/scaled2LariosCutoffFwdResults/seminorm_vs_t.png"}
\caption{$\|\nabla\ua\|_{L^2}$ vs $t$.}
\end{figure}


\section{Evolution of $\|\omega\|_{L^\infty}$}

This plot shows the evolution of the $\|\omega\|_{L^\infty}$ norm over time for each value of $\alpha$. This plot is compared to Fig. 1b in Bustamente et al. 2012. Additionally, we include our reproduction of Fig 1b for Euler. The plot shows that the evolution of the vorticity for the second scaled Euler-Voigt model does not converge to the plot in Fig 1b. Therefore, when the initial TGV is scaled to $\|\bn\eta_{TGV}\|_{L^2} = \pi$ the solution to Euler-Voigt does not match the results of Bustamante et al. 2012.

\begin{figure}[htbp]
    \centering
    \begin{subfigure}{0.7\textwidth}
        \includegraphics[width=\linewidth]{"./Figures/scaled2LariosCutoffFwdResults/vorticity_vs_t.png"}
        \caption{Scaled Euler-Voigt $\|\omega\|_{L^\infty}$ vs $t$.}
    \end{subfigure}
    \vspace{0.5cm} % Vertical space
    \begin{subfigure}{0.6\textwidth}
        \includegraphics[width=\linewidth]{"./Figures/Figure_2.png"}
		\caption{Scaled Euler $\|\omega\|_{L^\infty}$ vs $t$. Resolutions: 128, 256, 512, 1024}
    \end{subfigure}
    \vspace{0.5cm} % Vertical space
    \begin{subfigure}{0.6\textwidth}
        \includegraphics[width=\linewidth]{"./Figures/interplay_vorticity_b.png"}
		\caption{Bustamante Fig 1b. Euler $\|\omega\|_{L^\infty}$ vs $t$. Resolutions: 512, 1024, 2048, 4096.}
    \end{subfigure}
    \caption{Comparison of vorticity}
\end{figure}



\pagebreak
\section{Reproduction of Fig 2}

This plot shows the conservation of the $\alpha-$energy for the Euler-Voigt model and the conservation of $\alpha-$energy from Fig. 2 of Larios et al. 2018.

\begin{figure}[h!]
\centering
\begin{subfigure}{.33\textwidth}
\includegraphics[scale=0.3]{"./Figures/scaled2LariosCutoffFwdResults/alpha_energy_16.png"}
\end{subfigure}%
\vspace{0.5cm}
\begin{subfigure}{.33\textwidth}
\includegraphics[scale=0.3]{"./Figures/scaled2LariosCutoffFwdResults/alpha_energy_8.png"}
\end{subfigure}%
\vspace{0.5cm}
\begin{subfigure}{1.0\textwidth}
\centering
\includegraphics[scale=0.75]{"./Figures/Larios_Figure_2.png"}
\end{subfigure}
\end{figure}

The plot shows that both the size of the $\alpha-$Energy and the time-scale it changes over do not match the behaviour of the model from the Larios et al. plot. The plot for $\alpha=1/256$ is not included because it became under-resolved.

\pagebreak
\section{Reproduction of Fig 3}

This plot shows our reproduction of Fig 3. from Larios et al. 2018. It shows that our results match. The scale of the y axis is the same as in our plot. The slope of the black line in the Larios plot is $-1$, the slope of ours is also $-1$.
\\~\\
\begin{figure}[h!]
\centering
\begin{subfigure}{0.6\textwidth}
\centering
\includegraphics[scale=0.55]{"./Figures/scaled2LariosCutoffFwdResults/Fig3.png"}
\caption{Scaled Euler-Voigt. Green dots indicate the model became under-resolved.}
\end{subfigure}%
\begin{subfigure}{0.4\textwidth}
\centering
\includegraphics[scale=0.5]{"./Figures/Fig3_larios.png"}
\caption{Larios et al. plot}
\end{subfigure}
\end{figure}



\section{Reproduction of Fig 3 Inset}

This plot shows our reproduction of Fig 3. inset from Larios et al. 2018. It shows that our results match. The scale of the y axis is the same. The minimum slope in the Larios plot is below $-1$, the minimum slope of ours is also below $-1$. Finally, the time at which the slope reaches its minimum value is approximately the same as the Larios model.
\\~\\
\begin{figure}[h!]
\centering
\includegraphics[scale=0.75]{"./Figures/scaled2LariosCutoffFwdResults/Fig3_inset.png"}
\end{figure}


\section{Conclusion}

The results of our second scaled Euler-Voigt model show that it does not match the time-scale of the Bustamante et al. model, or the results of Fig 2, but does match the results of Larios et a. Fig 3.